\subsection{Semi-continuity of the volume of the homogeneous space}

In this No., for every measure \(\alpha \in \Gamma\) we set

\[
\mathrm{Q} _ {\alpha} = \mathrm{G} / \mathrm{H} _ {\alpha},\tag{1}
\]

and we denote by \(\pi_{\alpha}\) the canonical mapping \(\mathbf{G} \to \mathbf{Q}_{\alpha}\) .

Let \(\Gamma^0\) be the subset of \(\Gamma\) formed by the measures \(\alpha\) such that the subgroup \(\mathbf{H}_{\alpha}\) of \(\mathbf{G}\) is unimodular; the elements of \(\Gamma^0\) are characterized by the fact that \(\alpha(f) = \alpha(\check{f})\) for every function \(f \in \mathcal{K}(\mathbf{G})\) (every function of \(\mathcal{K}(\mathbf{H}_{\alpha})\) being extendible to a function of \(\mathcal{K}(\mathbf{G})\) by Urysohn's theorem); it follows that \(\Gamma^0\) is a closed subset of \(\Gamma\) . Recall that for every \(\alpha \in \Gamma^0\) , the quotient measure \(\mu_{\alpha} = \mu / \alpha\) on \(\mathbf{Q}_{\alpha}\) is defined and is relatively invariant under \(\mathbf{G}\) (Ch. VII, §2, No.~6, Th. 3); also recall that for every function \(f \in \mathcal{K}(\mathbf{G})\) ,

\[
\int_ {\mathbf {G}} f (x) d \mu (x) = \int_ {\mathbf {Q} _ {\alpha}} d \mu_ {\alpha} (\dot {x}) \int_ {\mathrm{H} _ {\alpha}} f (x s) d \alpha (s),\tag{2}
\]

where \(\dot{x} = \pi_{\alpha}(x)\) is the canonical image of \(x\in \mathbf{G}\) in \(\mathbf{Q}_{\alpha}\) .

PROPOSITION 4. --- Let \(\Gamma^0\) be the set of measures \(\alpha \in \Gamma\) such that \(\mathrm{H}_{\alpha}\) is unimodular, and for every \(\alpha \in \Gamma^0\) set \(\mu_{\alpha} = \mu / \alpha\) ; then the mapping \(\alpha \mapsto \| \mu_{\alpha} \|\) of \(\Gamma^0\) into \(\overline{\mathbf{R}}\) is lower semi-continuous for the vague topology.

For every \(\alpha \in \Gamma^0\) and every function \(f\in \mathcal{K}(\mathbf{G})\) , set

\[
f _ {\alpha} (\dot {x}) = \int_ {\mathrm{H} _ {\alpha}} f (x s) d \alpha (s) = (f * \alpha) (x),
\]

where the convolution product is taken relative to the right Haar measure \(\mu\) and where one makes use of the fact that \(\check{\alpha} = \alpha\) (§4, No.~4, formula (11)). We know (Ch. VII, §2, No.~1, Prop. 2) that the mapping \(f \mapsto f_{\alpha}\) of \(\mathcal{H}_{+}(G)\) into \(\mathcal{H}_{+}(Q_{\alpha})\) is surjective; therefore, by (2),

\[
\| \mu_ {\alpha} \| = \sup _ {f \in \mathcal {K} _ {+} (\mathrm{G}), f \neq 0} \mu_ {\alpha} (f _ {\alpha}) / \| f _ {\alpha} \| = \sup _ {f \in \mathcal {K} _ {+} (\mathrm{G}), f \neq 0} \mu (f) / \| f _ {\alpha} \|,
\]

where one has set

\[
\left\| f _ {\alpha} \right\| = \sup _ {\dot {x} \in \mathrm{Q} _ {\alpha}} | f _ {\alpha} (\dot {x}) | = \sup _ {x \in \mathrm{G}} | (f * \alpha) (x) |.\tag{3}
\]

To establish the proposition, it will suffice to show that, given \(f \in \mathcal{K}_{+}(G)\) , the mapping \(\alpha \mapsto \|f_{\alpha}\|\) is vaguely continuous. Now, let K be the support of f; the function \(f * \alpha\) has its support contained in \(KH_{\alpha}\) and is invariant on the right under \(H_{\alpha}\) ; consequently

\[
\left\| f _ {\alpha} \right\| = \sup _ {x \in K} | (f * \alpha) (x) |.
\]

The conclusion therefore follows from the fact that the mapping \(\alpha\mapsto f*\alpha\) of \(\mathcal{M}_{+}(\mathrm{G})\) equipped with the vague topology, into \(\mathcal{C}(\mathrm{G})\) equipped with the topology of compact convergence, is continuous (§4, No.~2, Remark 1).

Recall that if, for a measure \(\alpha \in \Gamma^0\) , \(\| \mu_{\alpha}\|\) is finite, then G is necessarily unimodular (Ch. VII, §2, No.~6, Cor. 3 of Th. 3).

PROPOSITION 5. --- Let g be a \(\mu\) -integrable positive numerical function and let \(\Gamma^{0}(g)\) be the set of measures \(\alpha\in\Gamma^{0}\) such that \(\int^{*}g(xs)d\alpha(s)\geqslant1\) for all \(x\in G\) . Then the mapping \(\alpha\mapsto\|\mu_{\alpha}\|\) of \(\Gamma^{0}(g)\) into \(\overline{\mathbf{R}}\) is vaguely continuous.

For every measure \(\alpha \in \Gamma^0 (\mathbf{G})\) , recall (Ch. VII, §2, No.~3, Prop. 5) that the function

\[
g _ {\alpha} (\dot {x}) = \int_ {\mathrm{H} _ {\alpha}} g (x s) d \alpha (s)
\]

is defined \(\mu_{\alpha}\) -almost everywhere on \(Q_{\alpha}\) , is \(\mu_{\alpha}\) -integrable, and

\[
\int_ {G} g (x) d \mu (x) = \int_ {Q _ {\alpha}} g _ {\alpha} (\dot {x}) d \mu_ {\alpha} (\dot {x}).\tag{4}
\]

In view of Prop. 4, it suffices to prove that, in \(\Gamma^0(g)\) , \(\alpha \mapsto \| \mu_\alpha \|\) is upper semi-continuous. Fix a measure \(\alpha \in \Gamma^0(g)\) , and let \(K\) be a compact subset of \(G\) . There exists on \(Q_\alpha\) a continuous function with compact support, taking its values in {[}0,1{]}, equal to 1 on the compact set \(\pi_\alpha(K)\) ; since the mapping \(f \mapsto f_\alpha\) of \(\mathcal{H}_+(\mathrm{G})\) into \(\mathcal{H}_+(\mathrm{Q}_\alpha)\) is surjective (Ch. VII, §2, No.~1, Prop. 2), one sees that there exists a function \(f \in \mathcal{H}_+(\mathrm{G})\) such that

\[
(f * \alpha) (x) = \int_ {\mathrm{G}} f (x s)   d \alpha (s) \left\{ \begin{array}{l} \leqslant 1 \text {for all} x \in \mathrm{G} \\ = 1 \text {for all} x \in \mathrm{K}. \end{array} \right.
\]

Since \(\beta \mapsto f * \beta\) is a continuous mapping of \(\mathcal{M}_{+}(\mathbf{G})\) , equipped with the vague topology, into \(\mathcal{C}(\mathbf{G})\) equipped with the topology of compact convergence (§4, No.~2, Remark 1), one sees that for every \(\varepsilon > 0\) , the set \(\mathrm{U}_{\varepsilon}\) of \(\beta \in \Gamma^0 (\mathbf{G})\) such that

\[
f _ {\beta} (\dot {x}) = \int_ {\mathrm{G}} f (x s) d \beta (s) > 1 - \varepsilon \quad \text { for   all } x \in \mathrm{K}
\]

is an open neighborhood of \(\alpha\) in \(\Gamma^0 (g)\) ; for every \(\beta \in \mathrm{U}_{\varepsilon}\) , we then have, by virtue of the formula (2),

\[
\left\| \mu_ {\alpha} \right\| \geqslant \int_ {G} f (x) d \mu (x) = \int_ {Q _ {\beta}} f _ {\beta} (\dot {x}) d \mu_ {\beta} (\dot {x}) \geqslant (1 - \varepsilon) \mu_ {\beta} \left(\pi_ {\beta} (K)\right).\tag{5}
\]

Given a number \(\varepsilon > 0\) , let us choose a function \(h \in \mathcal{H}_{+}(\mathrm{G})\) such that \(\int_{\mathrm{G}} |g(x) - h(x)| d\mu(x) \leqslant \varepsilon\) , and let us take \(\mathrm{K} = \operatorname{Supp}(h)\) in the foregoing. For every \(\beta \in \Gamma^0(g)\) , by hypothesis \(g_\beta(\dot{x}) \geqslant 1\) almost everywhere (for \(\mu_\beta\) ) in \(\mathbf{Q}_\beta\) , therefore

\[
\mu_ {\beta} \big (\mathrm{Q} _ {\beta} - \pi_ {\beta} (\mathrm{K}) \big) \leqslant \int_ {\mathrm{Q} _ {\beta} - \pi_ {\beta} (\mathrm{K})} g _ {\beta} (\dot {x}) d \mu_ {\beta} (\dot {x}) = \int_ {\mathrm{G} - \mathrm{KH} _ {\beta}} g (x) d \mu (x)
\]

by virtue of (4); since \(h\) is zero outside \(\mathbf{K}\) , and a fortiori outside \(\mathrm{KH}_{\beta}\) , it follows that

\[
\begin{array}{c} \mu_ {\beta} \big (\mathrm{Q} _ {\beta} - \pi_ {\beta} (\mathrm{K}) \big) \leqslant \int_ {\mathrm{G} - \mathrm{KH} _ {\beta}} | g (x) - h (x) | d \mu (x) \\ \leqslant \int_ {\mathrm{G}} | g (x) - h (x) | d \mu (x) \leqslant \varepsilon ; \end{array}
\]

combining this result with (5), one sees that

\[
\| \mu_ {\beta} \| \leqslant \varepsilon + \| \mu_ {\alpha} \| / (1 - \varepsilon)
\]

when \(\beta \in \mathrm{U}_{\varepsilon}\) , which completes the proof.

COROLLARY 1. --- Let K be a compact subset of G, V a symmetric compact neighborhood of e in G, c a real number \textgreater{} 0. The restriction of the mapping \(\alpha \mapsto \|\mu_{\alpha}\|\) to the set of \(\alpha \in \Gamma^0\) such that \(G = KH_{\alpha}\) and \(\alpha(V) \geqslant c\) is vaguely continuous.

For, let \(g \in \mathcal{K}_{+}(\mathbf{G})\) be a function such that \(g(x) \geqslant 1 / c\) for \(x \in \mathrm{KV}\) . For every \(x \in \mathbf{K}\) ,

\[
\int g (x s) d \alpha (s) \geqslant \int_ {\mathrm{V}} g (x s) d \alpha (s) \geqslant 1
\]

for \(\alpha\) satisfying the conditions of the statement; since, moreover, \(\pi_{\alpha}(\mathbf{K}) = \mathrm{Q}_{\alpha}\) , one has \(\alpha \in \Gamma^0(g)\) , whence the corollary.

COROLLARY 2. --- Let A be a \(\mu\) -integrable subset of G. The restriction of the mapping \(\alpha \mapsto \|\mu_{\alpha}\|\) to the set \(N_A\) of normalized Haar measures of the discrete subgroups H of G such that \(G = AH\) , is vaguely continuous.

For \(a \in \mathbf{A}\) and \(\alpha \in \mathbf{N}_{\mathbf{A}}\) ,

\[
\int \varphi_ {\mathrm{A}} (a s) d \alpha (s) \geqslant \varphi_ {\mathrm{A}} (a) = 1,
\]

and since \(\pi_{\alpha}(\mathbf{A}) = \mathbf{Q}_{\alpha}\) , one has \(\mathrm{N}_{\mathrm{A}} \subset \Gamma^{0}(\varphi_{\mathrm{A}})\) , and the corollary therefore follows from Prop. 5.

